%\documentclass[landscape,a0b,final,a4resizeable]{a0poster}
% final
\documentclass[spanish,portrait,a0b,final]{a0poster}
% resize
%\documentclass[portrait,a0b,final,a4resizeable]{a0poster}
%\documentclass[portrait,a0b,final]{a0poster}
%%% Option "a4resizeable" makes it possible ot resize the
%   poster by the command: psresize -pa4 poster.ps poster-a4.ps
%   For final printing, please remove option "a4resizeable" !!

\usepackage{epsfig}
\usepackage[usenames]{color}
% \usepackage{mathabx}  % symbols
\usepackage{amsmath}
\usepackage{amssymb}  % symbols
\usepackage{bm}
\usepackage{multicol}
\usepackage{pstricks,pst-grad}
\usepackage[latin1]{inputenc} % Allows typing accents in the editor directly, without \'a 
\usepackage{url} % Allows typesetting of url directions 
%\usepackage{babel}            % Language goes in the \documentclass line 
%\usepackage{babel}{spanish}


%%%%%%%%%%%%%%%%%%%%%%%%%%%%%%%%%%%%%%%%%%%
% Definition of some variables and colors
%\renewcommand{\rho}{\varrho}
%\renewcommand{\phi}{\varphi}
\setlength{\columnsep}{1.5cm}
\setlength{\columnseprule}{2mm}
\setlength{\parindent}{0.0cm}


%%% ----------------- Parameters and macro definitions ---------------------

\textwidth = 88cm

% Title boxes
\def\ltitu{9.5cm} % Length of the title box
% Main columns
\def\lcaja{84.0cm} % Length of the main  boxes  

%%%%%%%%%%%%%%%%%%%%%%%%%%%%%%%%%%%%%%%%%%%%%%%%%%%%
%%%               Background                     %%%
%%%%%%%%%%%%%%%%%%%%%%%%%%%%%%%%%%%%%%%%%%%%%%%%%%%%

\newcommand{\background}[3]{
\newrgbcolor{cgradbegin}{#1}
\newrgbcolor{cgradend}{#2}
\psframe[fillstyle=gradient,gradend=cgradend,
gradbegin=cgradbegin,gradmidpoint=#3](0.,0.)(1.\textwidth,-1.\textheight)
}



%%%%%%%%%%%%%%%%%%%%%%%%%%%%%%%%%%%%%%%%%%%%%%%%%%%%
%%%                Poster                        %%%
%%%%%%%%%%%%%%%%%%%%%%%%%%%%%%%%%%%%%%%%%%%%%%%%%%%%

%%   \begin{minipage}[c]{0.98\textwidth} ori
\newenvironment{poster}{
\begin{center}
\begin{minipage}[c]{0.98\textwidth}
}{
\end{minipage} 
\end{center}
}

%%%%%%%%%%%%%%%%%%%%%%%%%%%%%%%%%%%%%%%%%%%%%%%%%%%%
%%%                pcolumn                       %%%
%%%%%%%%%%%%%%%%%%%%%%%%%%%%%%%%%%%%%%%%%%%%%%%%%%%%

\newenvironment{pcolumn}[1]{
\begin{minipage}{#1\textwidth}
\begin{center}
}{
\end{center}
\end{minipage}
}

%%%%%%%%%%%%%%%%%%%%%%%%%%%%%%%%%%%%%%%%%%%%%%%%%%%%
%%%                pbox                          %%%
%%%%%%%%%%%%%%%%%%%%%%%%%%%%%%%%%%%%%%%%%%%%%%%%%%%%

\newrgbcolor{lcolor}{0. 0. 0.80}
\newrgbcolor{gcolor1}{1. 1. 1.}
\newrgbcolor{gcolor2}{.80 .80 1.}

\newcommand{\pbox}[4]{
\psshadowbox[#3]{
\begin{minipage}[t][#2][t]{#1}
#4
\end{minipage}
}}



%%%%%%%%%%%%%%%%%%%%%%%%%%%%%%%%%%%%%%%%%%%%%%%%%%%%
%%%                myfig                         %%%
%%%%%%%%%%%%%%%%%%%%%%%%%%%%%%%%%%%%%%%%%%%%%%%%%%%%
% \myfig - replacement for \figure
% necessary, since in multicol-environment 
% \figure won't work

\newcommand{\myfig}[3][0]{
\begin{center}
\vspace{0.5cm}
\includegraphics[clip,width=#3\hsize,angle=#1]{#2}
\nobreak\medskip
\end{center}}



%%%%%%%%%%%%%%%%%%%%%%%%%%%%%%%%%%%%%%%%%%%%%%%%%%%%
%%%                mycaption                     %%%
%%%%%%%%%%%%%%%%%%%%%%%%%%%%%%%%%%%%%%%%%%%%%%%%%%%%
% \mycaption - replacement for \caption
% necessary, since in multicol-environment \figure and
% therefore \caption won't work

%\newcounter{figure}
\setcounter{figure}{1}
\newcommand{\mycaption}[1]{
  \vspace*{-1cm}
  \begin{quote}
%ori    {{\sc Figure} \arabic{figure}: #1}
    {\footnotesize #1}
  \end{quote}
  \vspace{0.5cm}
  \stepcounter{figure}
}



%%%%%%%%%%%%%%%%%%%%%%%%%%%%%%%%%%%%%%%%%%%%%%%%%%%%%%%%%%%%%%%%%%%%%%
%%% Begin of Document
%%%%%%%%%%%%%%%%%%%%%%%%%%%%%%%%%%%%%%%%%%%%%%%%%%%%%%%%%%%%%%%%%%%%%%

\begin{document}

\background{1. 1. 1.}{1. 1. 1.}{0.5}

%\vspace*{2cm}


\newrgbcolor{lightblue}{0. 0. 0.80}
\newrgbcolor{white}{1. 1. 1.}
\newrgbcolor{black}{0. 0. 0.}
\newrgbcolor{whiteblue}{.80 .80 1.}
\newrgbcolor{orange}{.10 .50 .50}


\begin{poster}
%
%%%%%%%%%%%%%%%%%%%%%  Header and Title %%%%%%%%%%%%%%%%%%%%
%
\begin{center} \begin{pcolumn}{0.98}
\pbox{0.95\textwidth}{\ltitu}{linewidth=2mm,framearc=0.3,linecolor=black,
fillstyle=gradient,gradangle=0,gradbegin=white,gradend=white,gradmidpoint=1.0,framesep=1em}{
%
%%% Left logo 
%
\begin{minipage}[c][9cm][c]{0.1\textwidth} \begin{center}
%\includegraphics[width=7.0cm,angle=0]{goe-logo.eps}
\vspace*{-2cm} 
\includegraphics[width=7.0cm,angle=0]{graphics/Logo_CONICET.eps}
\end{center} 
\end{minipage}
%%% Title
\begin{minipage}[c][9.5cm][c]{0.78\textwidth} 
\begin{center} 
\vspace*{0.0cm}
{\Huge \bf \sc Droplet transport in a polymer-coated nanochannel }\\[4mm] 
\vspace*{0.2cm}
{\Large  {Kevin Speyer}$^{*}$ and \underline{Claudio Pastorino}$^{*}$ \\[5mm] 
\vspace*{-0.5cm}
$^{*}$Departamento de Materia Condensada,  Centro At�mico Constituyentes CNEA-CONICET, Buenos Aires, Argentina \\[4mm]}
\end{center}
\vspace*{-0.5cm}
\hspace*{2cm}
pastor@cnea.gov.ar
\vspace*{-0.5cm}

%%} % small

\end{minipage}
%
%%% Right logo 
%
\begin{minipage}[c][12cm][c]{0.1\textwidth} \begin{center}

\vspace*{-2.5cm} \hspace*{1.0cm} \includegraphics[width=5.0cm,angle=0]{graphics/Logo_CNEA.eps}

 \hspace*{1.0cm} \includegraphics[width=5.0cm,angle=0]{graphics/logo-MC_150.eps} \end{center}
\end{minipage}
} \end{pcolumn}
\end{center}
\vspace*{0.2cm}


%%%%%%%%%%% Ends poster header %%%%%%%%%%%%%%%%%%

%%%%%%%%%%%%%%%%%%%%%
%%% Content
%%%%%%%%%%%%%%%%%%%%%

%%%%%%%%%%%%%%%%%%%%%%%%%%%%%%% ABSTRACT  %%%%%%%%%%%%%%%%%%%%%%%%%%%%%%%%%
\vspace{0.0cm}
\begin{center}
\begin{pcolumn}{0.87}    % one column abstract 
\pbox{0.95\textwidth}{12cm}{linewidth=2mm,framearc=0.3,linecolor=blue,fillstyle=gradient,gradangle=0,gradbegin=white,gradend=white,gradmidpoint=1.0,framesep=1em}
{
\begin{center}
\vspace{0.3cm} 
{\Large \sl \bf Abstract } 
\end{center}
\vspace{0.7cm}
We study the influence of chain stiffness on droplet flow in a nano-channel,
coated with semiflexible hydrophobic polymers by means of non-equilibrium
molecular-dynamics simulations. The studied system is then a moving
droplet in the slit channel, coexisting with its vapor and subjected
to periodic boundary conditions in the flow direction. The polymer
chains, grafted by the terminal bead to the confining walls, are described
by a coarse-grained model that accounts for chain connectivity, excluded
volume interactions and local chain stiffness. The rheological, frictional
and dynamical properties of the brush are explored over a wide range
of persistence lengths. We find a rich behavior of polymer conformations
and concomitant changes in the friction properties over the wide range
of studied polymer stiffnesses. A rapid decrease in the droplet velocity
was observed as the rigidity of the chains is increased for polymers
whose persistence length is smaller than their contour length. We
find a strong relation between the internal dynamics of the brush
and the droplet transport properties, which could be used to tailor
flow properties by surface functionalization. The monomers of the
brush in the boundary layer under the droplet, present a collective ``treadmill
belt'' like dynamics which can only be present due the the existence
of grafted chains. We describe its changes in spatial extension upon
variations of polymer stiffness, with bidimensional velocity and density
profiles. The deformation of the polymer brushes due to the presence
of the droplet is analyzed in detail. Lastly, the droplet-gas interaction
is studied by varying the liquid to gas ratio, observing a 16\% speed
increase for droplets that flow close to each other, compared to a
train of droplets that present a large gap between consecutive droplets. 
}
\end{pcolumn}
\end{center}

\begin{pcolumn}{0.5}    % first coulmn of the poster 1 1
\pbox{0.9\textwidth}{\lcaja}{linewidth=2mm,framearc=0.1,linecolor=black,fillstyle=gradient,gradangle=0,gradbegin=white,gradend=white,gradmidpoint=1.0,framesep=1em}
{
%%% * **************************   Introduction   *****************************

        \vspace{0.0cm}
        \begin{center}
        \pbox{0.8\textwidth}{}{linewidth=2mm,framearc=0.1,linecolor=black,fillstyle=gradient,gradangle=0,gradbegin=white,gradend=white,gradmidpoint=1.0,framesep=1em}{
        \begin{center}
        {\LARGE \bf Introduction }
        \end{center}
        }
        \end{center}
        \vspace{1.0cm}

\begin{center}  {\bf \normalsize \textcolor{red}{\Large System of Interest}  } \end{center}

	            \hspace*{0cm}
          \begin{minipage}[t]{0.95\textwidth} %Left minipage start
          \vspace*{0cm}  
          \myfig[0]{graphics/Fig1_1.eps}{0.9}
          \end{minipage} \hspace*{0cm}  
          
Snapshot of the simulated nano-channel. The semiflexible polymer chains (in \textcolor{red}{ \bf red}) are grafted by
the end bead (in {\color{BlueGreen} \bf light blue}) to the wall of the channel. The fluid (droplet + vapor)
inside the channel is composed by particles (in {\color{gray} \bf gray}) interacting via a Lennard-Jones potential.
Periodic boundary conditions are imposed in directions $\hat{x}$ and $\hat{y}$. An external constant force is applied
to all particles in $\hat{x}$ direction.

%           
%          \begin{minipage}[t]{0.95\textwidth}
%          \begin{center}
%        \vspace*{0cm}
%            \myfig[0]{graphics/Brushes3.eps}{0.9}
%          \end{center}
%          \end{minipage} % Right minipage end
%
%Polymers of varying rigidities were used to investigate the friction between droplet and polymer brush.
%The bending constant $k_{\rm b}$ took 3 different values to explore different regimes of persistence
%length over contour length: $l_{\rm p}/l_{\rm c} \simeq 0.1; 1;10$. The interaction between polymers and fluid is
%hydrophobic.
%

        \vspace{2.0cm}
        \begin{center}
        \pbox{0.8\textwidth}{}{linewidth=2mm,framearc=0.1,linecolor=black,fillstyle=gradient,gradangle=0,gradbegin=white,gradend=white,gradmidpoint=1.0,framesep=1em}{
        \begin{center}
        {\LARGE \bf Model and Simulation Techniques }
        \end{center}
        }
        \end{center}

         % \begin{minipage}{1cm}
          %\vspace*{0.7cm}
         %{\Huge \textcolor{red}{ $\Rightarrow$} }
          %\end {minipage}
          \vspace{1.0cm} 
          \begin{minipage}{40cm}
          \begin{itemize}
          \item Molecular Description: Coarse grain, bead spring model. 
          \item \textcolor{blue}{\bf \normalsize Molecular Dynamics (MD) + Dissipative Particle Dynamics Thermostat (DPD)}\\[0.5cm]
          \end{itemize}
          \end{minipage}
        \vspace*{0cm}
    %     {\bf \textcolor{red}{ \bf Sistema modelo:} }

%         Consideramos el caso en que el solvente (l�quido polim�rico) y las cadenas
%         del cepillo tienen incompatibilidad qu�mica y forman una interfase en la cu�l el l�quido no entra completamente
%         en el interior del cepillo.
%         \framebox{\bf T�cnica de Simulaci�n:} \textcolor{blue}{ \bf \normalsize MD + termostato DPD }
%         para dar cuenta de las interacciones hidrodin�micas (conserva momento localmente)\cite{claudio2}.
\vspace*{-1cm}
\begin{center} {\bf \textcolor{red}{ \bf \Large Interactions} } \end{center}
        \begin{itemize}
         \item {\bf Chain Connectivity} between contiguous {\sl beads} of the same polymer.
             Finitely extensible non-linear elastic potential(FENE):

\vspace*{0cm}
\begin{minipage}[t]{0.48\textwidth} 
\vspace*{-1cm}

		\[ U_{\rm FENE}=  \begin{cases} -\frac{1}{2}k\ R_{0}^{2}\ln\left[1-(\frac{r}{R_{0}})^{2}\right] 
		& r\leq R_{0}\\ \infty & r>R_{0} \end{cases} , \] 
\end{minipage} \begin{minipage}[t]{0.48\textwidth}
\vspace*{1cm}
where $R_0=1.5\sigma$, $k=30\varepsilon/\sigma^2$ and $r=|{\mathbf r_{i}}-{\mathbf r{j}}|$ .
\end{minipage} 
       %Lennard-Jones 

\vspace*{1cm}
       \item {\bf Monomer-Monomer Potential}
           \vspace*{0cm}
        \hspace*{0cm} 

\begin{minipage}[t]{0.48\textwidth}
\vspace*{-1cm}
           \[ U(r) =  \begin{cases} 4\varepsilon[(\frac{\sigma}{r})^{12} -(\frac{\sigma}{r})^6 -U_{LJ}(r_c) ] & r\leq r_c
          \\ 0  & r > r_c  \end{cases},\]
%\vspace*{1.1cm} %\varepsilon^{l�q-l�q}= 1 \varepsilon^{cep-cep}$

          \end{minipage} 
%\hspace*{-1cm}   
\begin{minipage}[t]{0.48\textwidth}
%\hspace{1cm}
\hspace{1cm}{ \hspace{1cm} \bf \normalsize \textcolor{blue}{Hidrophobic}} liquid-brush interaction
\[\varepsilon^{liq-cep}=\frac{1}{3}\varepsilon^{liq-liq}=\frac{1}{3}\varepsilon^{cep-cep}\]
 
        %\begin{center}
%        \vspace*{-6cm}
         %   \myfig[0]{graphics/Potenciales.eps}{0.85}
        %  \end{center}
 \end{minipage} % Right minipage
        \vspace*{0.4cm}
           \\ $\sigma=1$. Polymer cut radius: $r_c=2^{\frac{1}{6}}$ (polymers in good solvent conditions).

          %Potencial de Flexi�n 
         
           \vspace*{1cm}

           \item {\bf Wall-Monomer Potential}

           \vspace*{0cm}
\hspace*{1.5cm} 
\begin{minipage}[t]{0.3\textwidth}
\vspace*{-0.5cm}  
\hspace*{0cm}         
\[ U_{\rm wall}(z) =  |A_w|(\frac{\sigma_w}{z})^{9} -A_w(\frac{\sigma_w}{z})^3 ,  \]
 \end{minipage} \hspace*{1cm}   \begin{minipage}[t]{0.65\textwidth}
\vspace*{0.5cm}          
 with  $A=-3.2\varepsilon$. The walls, fixed at $z=0$ and $z=D$,

  are impenetrable and have no intrinsic roughness.
	\end{minipage} % Right minipage
\vspace*{0.5cm}
\item {\bf Bending Potential}

            \begin{minipage}[t]{0.68\textwidth}
            \begin{center}
            \vspace*{-1.5cm}
             \[ U_{\rm bend}(\theta) =  \frac{1}{2}  k_{\rm b} \theta^{2} \]
            \end{center}
        This potential favours chain stretching and adds rigidity to the polymer. A wide range of bending
	constants ($k_{\rm b}$) were implemented to explore the different rigidity regimes:
        $ 0 < k_{\rm b} < 160\varepsilon $. In terms of persistence length ($l_p$) over contour length ($l_c$):
        $ 0 < l_p/l_c < 22 $.

        The reduced persistence length ($l_p/l_c$) is the dimensionless quantity that meassures the
        polymer rigidity. 
        %    \end{center}
          \end{minipage} \hspace*{1cm}   \begin{minipage}[t]{0.3\textwidth}
          \begin{center}
        \vspace*{-0.0cm}
            \myfig[0]{graphics/Esquema_Bending.eps}{0.6}
          \end{center}
          \end{minipage} % Right minipage

           \end{itemize}
         \vspace*{0.cm} 

%_TABLA_%
%\begin{center} 
%\vspace{0cm}
%    \begin{tabular}{ | c | c | c | c | c |}
%    \hline
%    $k_b$ [$\varepsilon$] & $k_b$/$k_BT$ & $l_p/l_c$  & $l_p/a$ & $\tau_R$ [$\sigma\sqrt{m/\varepsilon}$] \\ \hline
%    0   & 0   & 0.16 & 1.4 & 17 \\ \hline    
%    8   & 10  & 1.12 & 10.1 & 50 \\ \hline
%    80  & 100 & 10.1 & 91 & 260 \\ \hline
%    \end{tabular}
%
%
%\vspace{1cm}
%Table: polymer properties as a function of the bending constant $k_b$: persistence length ($l_p$) and relaxation time ($\tau_R$). Chain contour length $l_c=8.55\sigma$ ($N=10$). Constant simulation temperature $k_bT=0.8\varepsilon$.
%\end{center}


}
\end{pcolumn}   \begin{pcolumn}{0.5} %%%%%%%%%%%%%%%%%%%%% Third column
\pbox{0.9\textwidth}{\lcaja}{linewidth=2mm,framearc=0.1,linecolor=black,fillstyle=gradient,gradangle=0,gradbegin=white,gradend=white,gradmidpoint=1.0,framesep=1em}{

\vspace{0cm}\begin{center}\pbox{0.8\textwidth}{}{linewidth=2mm,framearc=0.1,linecolor=black,fillstyle=gradient,gradangle=0,gradbegin=white,gradend=white,gradmidpoint=1.0,framesep=1em}{\begin{center}{\LARGE \bf Results}\end{center}}\end{center}\vspace{0.25cm}

      
\vspace*{1cm}

%Fluid Velocity Field
\begin{center} {\bf \Large \textcolor{red}{Polymer Dynamics}} \end{center}

% \begin{minipage}[t]{0.23\textwidth}
 \vspace*{0cm}
% \myfig[0]{./graphics/Rend3.eps}{0.95}
%\end{minipage}  \hspace*{2cm}
\begin{center}
 \begin{minipage}[t]{0.99\textwidth}
	% \vspace*{0cm}
\myfig[0]{./graphics/Rend_p_histo_2D_3.eps}{0.99}
\end{minipage}
\end{center}

%Epigraph of Image
Polymer's end-bead distribution is shown
as a color map. The correspondence between colors and monomer density
(given in units of $\sigma^{-3}$) is shown in the color bar at the
bottom. The black arrows, composing a vector field, represent the mean
momentum of the end-beads of the polymers. These graphs are presented
for various rigidities, from left to right and top to bottom: $l_{p}/l_{c}=0.1;\,0.72;\,2.9;\,8.3$.


\vspace*{1.5cm}
% DOPLET VELOCITY VS L
\begin{center} {\bf \Large \textcolor{red}{Droplet Velocity vs Polymer Rigidity}}\end{center}

\begin{center}

\begin{minipage}[t]{0.45\textwidth}
% \vspace*{0.4cm}
\hspace{0cm}
\myfig[0]{./graphics/dropVel_vs_rigidity.eps}{0.99}
Droplet velocity with respect
to the reduced persistence length of the polymers. The inset shows
the droplet velocity plotted against the persistence length in a semi-log
graph to highlight the behavior at low persistence lengths. For $l_p/l_c<2$ the droplet
velocity decreases rapidly with increasing polymer rigidity. The liquid velocity is highly
influenced by the polymer dynamics. 
\end{minipage}
\hspace{2cm}
\begin{minipage}[t]{0.45\textwidth}
% \vspace*{0.4cm}
\hspace{0cm}
\myfig[0]{./graphics/SlipVel_vs_rigidity3.eps}{0.99}
Velocity of the brush (black circles)
and liquid (red squares) in the brush-liquid interface are plotted
with respect to the reduced persistence length. The difference between
these velocities, shows how the liquid boundary layer passes over
the brush boundary layer (open green diamonds). Inset: droplet velocity
versus the liquid boundary layer velocity.
\end{minipage}
\end{center}
%Friction Forces FIG
\vspace*{1.5cm}
 
\begin{center} {\bf \Large \textcolor{red}{Flow Across the Channel}}\end{center}

\begin{center} 
%\hspace*{0cm}
\vspace*{0.0cm} 
\begin{minipage}[t]{0.42\textwidth}
 \vspace*{1.0cm}
Fluid velocity in the direction of the
flow ($\hat{x}$) as a function of the x coordinate (in log scale).
The droplet ($x<48\sigma$) travels at a higher velocity than the
vapor phase ($x>48\sigma$). It is also possible to observe that the gas
velocity decays exponentially with x. Inset: 2-dimensional density profile of the droplet, with the liquid
velocity field overlayed. The velocities are calculated in the frame
of reference of the center of mass of the droplet. The color legend
shows the liquid number density in units of $\sigma^{-3}$.

\end{minipage} 
\hspace*{0.5cm} 
%\vspace*{0.0cm}
\begin{minipage}[t]{0.55\textwidth}
\hspace*{1.6cm}

\vspace*{-2.1cm}
\myfig[0]{./graphics/vfluid_vs_Lx_4.eps}{0.99}
\end{minipage}
\end{center}

\vspace*{0cm}


%%% %%%%%%%%%%%%%%%%% BIB %%%%%%%%%%%%%%
\vspace*{-1cm}
\textbf{\Large References}
 
K. Speyer and C. Pastorino, \textit{Soft Matter}, 2015, DOI: 10.1039/C5SM01075F

K. Speyer and C. Pastorino, \textit{Langmuir}, 2017, DOI: 10.1021/acs.langmuir.7b02640
 \vspace*{-1.3cm}
%\renewcommand{\refname}{Referencias}
%{\small
%\begin{thebibliography}{1}
%
%\bibitem{0295-5075-102-5-58003}
%Andrey Milchev,  Kurt Binder.
%\newblock {\em EPL (Europhysics Letters)}, 102(5):58003, 2013.
%
%\bibitem{0953-8984-23-18-184105}
%F~Leonforte, J~Servantie, C~Pastorino, M~M{\"u}ller.
%\newblock {\em Journal of Physics: Condensed Matter}, 23(18):184105, 2011.
%
%\bibitem{:/content/aip/journal/jcp/140/1/10.1063/1.4851195}
%C.~Pastorino,  M.~M{\"u}ller.
%\newblock {\em The Journal of Chemical Physics}, 140(1):--, 2014.
%
%\bibitem{pastorino:064902}
%C.~Pastorino, K.~Binder, T.~Kreer, M.~M{\"u}ller.
%\newblock {\em The Journal of Chemical Physics}, 124(6):064902, 2006.
%
%\bibitem{PhysRevE.76.026706}
%C.~Pastorino, T.~Kreer, M.~M\"uller, K.~Binder.
%\newblock {\em Phys. Rev. E}, 76:026706, Aug 2007.
%
%\end{thebibliography}
%
%
%
%
%% \begin{thebibliography}{77}
%% 
%%%\bibitem{claudio3} M.  M\"uller and C. Pastorino,
%%% {\sl Europhys. Lett. }  {\bf 81}, 28002 (2008).
%%
%%\bibitem{claudio2} C. Pastorino, T. Kreer, M.  M\"uller and K. Binder, 
%% {\sl Phys. Rev. E}  {\bf 76}, 026706 (2007).
%%
%%\bibitem{claudio1} C. Pastorino, K. Binder, T. Kreer, and M.
%%M\"uller, J. Chem. Phys. {\bf 124},  64902 (2006)   
%%
%%\bibitem{claudio4} C. Pastorino, K. Binder, and M.  M\"uller, 
%%{\bf submitted to } {\sl Macromolecules} {\bf 42}, 401 (2009)
%%\bibitem{fabien} F. L�onforte, J. Servantie, C. Pastorino, and M. M�ller ,
%%aceptado en  {\sl J. Phys. Cond. Matt.} (2010)
%%\bibitem{doyle} P. S. Doyle, B. Ladoux and J. L. Viovy,  {\sl Phys. Rev. Lett. } {\bf 84}, 4769 (2000)
%%%\bibitem{winkler} R. G. Winkler , {\sl Phys. Rev. Lett.} {\bf 97}, 128301 (2006)
%%\bibitem{gerashenko} S. Gerashenko and V. Steinberg, {\sl Phys. Rev. Lett.}  {\bf 96}, 38304 (2006)
%% \end{thebibliography}
%} %end small  
 }
\end{pcolumn}
%  ********************************** END OF POSTER *********************************** 
\end{poster}
% *************************************** END  ***************************************
\end{document}

